\documentclass[a4paper, 12pt]{report}

\usepackage{charter}
\usepackage{makeidx}
\usepackage{fancyhdr}
\usepackage{hyperref}
\usepackage[utf8]{inputenc}
\usepackage{graphicx}
\usepackage[left=2cm, right=2cm]{geometry}
\usepackage{latexsym}
\usepackage{amsmath, amsthm, amssymb}
\usepackage{rotating}


\begin{titlepage}
\title{Requisiti del Prodotto}
\author{Release 1.0}
\date{\today \\Firenze \\\begin{figure}[h] \centering
\includegraphics[width=0.2\textwidth]{../images/logokiwi.png} \end{figure} }
\end{titlepage}

\pagestyle{fancy}

\begin{document}

\maketitle

\section*{Approvazione, redazione, lista distribuzione}
\begin{table}[h!]
  \begin{center}
    \begin{tabular}{| l | l | p{60mm} |}
    \hline
    \textbf{approvato da} & \textbf{il giorno} & \textbf{firma} \\
	\hline    
	Marco Tinacci & 16/09/2009 &  \\
    \hline
    \end{tabular}
  \end{center}
\end{table}

\begin{table}[h!]
  \begin{center}
    \begin{tabular}{| l | l | p{60mm} |}
    \hline
    \textbf{redatto da} & \textbf{il giorno} & \textbf{firma} \\
    \hline
    Francesco Calabri & 16/09/2009 &  \\
    \hline
	Manuele Paulantonio & 16/09/2009 &  \\
    \hline    
	Massimo Nocentini & 16/09/2009 &  \\
    \hline
    \end{tabular}
  \end{center}
\end{table}

\begin{table}[h!]
  \begin{center}
    \begin{tabular}{| l | l | p{60mm} |}
    \hline
    \textbf{distribuito a} & \textbf{il giorno} & \textbf{firma} \\
	\hline    
	Daniele Poggi & 16/09/2009 &  \\
    \hline
	Niccol\'o Rogai & 16/09/2009 &  \\
    \hline
	Marco Tinacci & 16/09/2009 &  \\
    \hline
    \end{tabular}
  \end{center}
\end{table}

\tableofcontents

\newpage

\section*{Introduzione}

\chapter{Requisiti obbligatori}

\section{Generali (2.1)}
Come requisiti fondamentali, \textbf{PMango 3.0} sar\`a visualizzabile e
usabile con le ultime versioni \emph{Internet Explorer 8} e \emph{Mozilla
Firefox 3.0}.

Le nostre modifiche e aggiunte saranno distribuite senza costi di licenza, in
quanto si tratta di estensioni di un progetto GPL

Ci assumiamo la responsabilit\`a di essere conformi ai punti \emph{d), e)}.

\section{Diagrammi WBS, Gantt e Task Network (2.2)}
\begin{itemize}
  \item[a)] implementato nello use case \textbf{1.3 Show Project Page}.
  \item[b)] implementato negli use case \textbf{1.2 Make NodeTaskbox} e in
  \textbf{2.2 Make GanttTaskbox}.
  \item[c)] implementato nello use case \textbf{1.5 Show UserOptions}
\end{itemize}

\section{Diagrammi specifici (2.3)}
\begin{itemize}
  \item[a)] implementato nello use case \textbf{1.5\footnote{replace this
  number with the relative use case within the tasknetwork package} Show
  UserOption}
  \item[b)] implementato nello use case \textbf{4.2 Create Critical Path Table}
  \item[c)] implementato nello use case \textbf{1.5 Show UserOptions}
  \item[d)] implementato nello use case \textbf{1.5 Show UserOptions} e
  descritto in modo dettagliato nella sezione \textbf{7.1 UserOption’s
  Instances} del documento \textbf{Domain Model}
  \item[e)] implementato nello use case \textbf{1.5 Show UserOptions}
\end{itemize}

\section{Generazione di immagini e doc (2.4)}
\begin{itemize}
  \item[a)] implementato negli use case \textbf{1.10 Refresh Chart, 1.9 Make
  PDF}
  \item[b)] implementato nello use case \textbf{1.8 Add to Report UserAction}
  \item[c)] implementato nello use case \textbf{1.3 Show UserOptions} e
  descritto in modo dettagliato nella sezione \textbf{7.1 UserOption’s
  Instances} del documento \textbf{Domain Model}
  \item[d)] vedi punto \emph{a)}
  \item[e)] implementato nello use case \textbf{1.12 Open in New Window}
\end{itemize}

\section{Documentazione (2.5)}
Ci assumiamo la responsabilit\`a di essere coerenti a quanto richiesto nei
punti \emph{a), b), c)} nei momenti in cui verranno effettivamente implementati.

\chapter{Semplificazioni, Metriche}
\section{Semplificazioni e requisiti aggiuntivi}
\begin{itemize}
  \item[a)] il nostro gruppo \textbf{non} prevede lo sviluppo di requisiti
  aggiuntivi, preferendo implementare correttamente il processo di sviluppo adottato per
  raggiungere i requisiti richiesti dal committente.
  \item[b)] abbiamo deciso di creare \textbf{solo} oggetti \textbf{gif} per
  usarli in modo interscambiabile sia nella visualizzazione da browser web, sia
  per aggiungerli in documenti PDF. Questo ci porta alcuni vantaggi:
  \begin{itemize}
    \item ci interfacciamo con una sola libreria, avendo cosi modo di capirne a
    fondo il comportamento e eventualmente aggiungere quelle funzionalit\`a di
    helper che potrebbero servirci, ma che attualmente non vengono fornite.
    \item ci riduce il carico di lavoro, questo non preclude che se arriviamo
    in anticipo con un prodotto finito e che rispetta la specifica richiesta,
    potremo proporre una integrazione dell'offerta sviluppando le
    funzionalit\`a native per la rappresentazione in PDF.
    \end{itemize}
\end{itemize}

\section{Metriche}

\end{document}